% =====================================================================
\chapter{Tensores, \emph{broadcasting} y formas de un LLM}
\label{tema:6}
% =====================================================================

En el código de un modelo de lenguaje casi nunca aparecen vectores o matrices aislados, sino arrays de varios ejes que agrupan muchos vectores a la vez. Este tema formaliza esos arrays, el producto matricial por lotes y las reglas de \emph{broadcasting}, y las aplica a las formas de las matrices de un Transformer y al rango de su capa de salida. Se sigue a Zhang \emph{et al.}~\cite[§2.1 y §2.3]{d2l}, la documentación de NumPy~\cite{numpy-bc} y el apéndice~6 de~\cite{strang-ila}; la última sección, a Yang \emph{et al.}~\cite{yang-softmax}.

\section{Tensores como arrays}

\begin{definicion}[Tensor]
Un \emph{tensor} de orden $k$ y forma $(n_1,\dots,n_k)$ es una familia de números reales $X[i_1,\dots,i_k]$ indexada por $0\le i_j<n_j$; contiene $n_1n_2\cdots n_k$ elementos. Cada índice recorre un \emph{eje}. Los escalares, vectores y matrices son tensores de orden 0, 1 y~2.
\end{definicion}

En física y en geometría diferencial, un tensor es una aplicación multilineal cuyas componentes cambian de una forma determinada al cambiar de base~\cite[apéndice~6]{strang-ila}. En aprendizaje automático el término designa simplemente un array multidimensional~\cite[§2.3.4]{d2l}, y es el sentido que se usa en este curso.

\begin{definicion}[Reducción a lo largo de un eje]
Sumar un tensor de forma $(n_1,\dots,n_k)$ a lo largo del eje $j$ produce un tensor de forma $(n_1,\dots,n_{j-1},n_{j+1},\dots,n_k)$:
\[
Y[i_1,\dots,\widehat{i_j},\dots,i_k]=\sum_{i_j=0}^{n_j-1}X[i_1,\dots,i_k],
\]
donde $\widehat{i_j}$ indica que el índice se suprime. Con la opción \texttt{keepdims=True}, el eje $j$ se conserva con tamaño~1.
\end{definicion}

La media, el máximo y la función $\LSE$ del módulo~1 se reducen del mismo modo. En un LLM, la softmax se aplica a lo largo del último eje de los logits y la normalización por capas calcula media y varianza a lo largo del eje de $\dmodel$.

\section{Producto por lotes y notación de Einstein}

\begin{definicion}[Producto matricial por lotes]
Sean $X$ de forma $(\dots,m,p)$ e $Y$ de forma $(\dots,p,q)$, con los mismos ejes iniciales (o ejes compatibles por \emph{broadcasting}). El producto $X\,@\,Y$ tiene forma $(\dots,m,q)$ y se obtiene multiplicando como matrices los dos últimos ejes para cada valor de los índices iniciales, que actúan como índices de lote.
\end{definicion}

Multiplicar una matriz de $m\times p$ por otra de $p\times q$ requiere $mpq$ multiplicaciones. Por eso el orden en que se agrupa un producto de varias matrices afecta mucho al coste (\cref{eje:6-orden}).

La notación de Einstein (\texttt{np.einsum}) describe estas operaciones nombrando los índices: los que aparecen en las entradas y no en la salida se suman. Por ejemplo, para $Q$ y $K$ de forma $(B,T,d)$, la expresión \texttt{einsum('bid,bjd->bij', Q, K)} calcula
\[
S[b,i,j]=\sum_{k}Q[b,i,k]\,K[b,j,k],
\]
es decir, $S=QK\T$ para cada elemento del lote, de forma $(B,T,T)$. Es el núcleo de la atención del módulo~13. Con \texttt{@}, el mismo cálculo exige trasponer solo los dos últimos ejes: \texttt{Q @ K.transpose(0, 2, 1)}.

\section{Broadcasting}

\begin{definicion}[Reglas de \emph{broadcasting}]\label{def:broadcasting}
Dos tensores de formas $(a_1,\dots,a_p)$ y $(b_1,\dots,b_q)$ se combinan elemento a elemento del siguiente modo~\cite{numpy-bc}:
\begin{enumerate}
  \item se alinean las formas por la derecha y se completa la más corta con ejes de tamaño~1 por la izquierda;
  \item en cada eje, los tamaños son \emph{compatibles} si son iguales o si uno de ellos vale~1; en otro caso la operación produce un error;
  \item el resultado tiene, en cada eje, el mayor de los dos tamaños, y el tensor con tamaño~1 en ese eje se repite a lo largo de él.
\end{enumerate}
\end{definicion}

\begin{table}[h]
\centering\small
\caption{Ejemplos de \emph{broadcasting} con activaciones de forma $(2,3,4)$.}\label{tab:broadcasting}
\begin{tabular}{@{}lll@{}}
\toprule
Operación & Formas & Resultado \\
\midrule
\texttt{x + b} & $(2,3,4)+(4,)$ & $(2,3,4)$: sesgo sumado a cada token \\
\texttt{x - x.mean(-1, keepdims=True)} & $(2,3,4)-(2,3,1)$ & $(2,3,4)$: centra cada token \\
\texttt{x @ W} & $(2,3,4)\,@\,(4,8)$ & $(2,3,8)$: misma capa en cada token \\
\texttt{x + y} & $(2,3,4)+(3,)$ & error: $4\ne 3$ \\
\bottomrule
\end{tabular}
\end{table}

\begin{observacion}
El \emph{broadcasting} es la forma eficiente de escribir operaciones que en notación matemática llevan un vector de unos: si $X\in\R^{T\times d}$ y $b\in\R^{d}$, \texttt{X + b} calcula $X+\mathbf{1}b\T$ sin construir la matriz $\mathbf{1}b\T$. Su contrapartida es que algunas combinaciones de formas no producen error y dan un resultado no deseado; el \cref{eje:6-formas} muestra un caso típico. Comprobar las formas en cada paso es el hábito que más errores evita al programar un Transformer.
\end{observacion}

\section{Formas de un LLM}

El \cref{tab:formas} recoge las formas de los tensores principales de un Transformer de tipo decodificador, con los valores de GPT-2 small~\cite{gpt2}: vocabulario de $V=50\,257$ tokens, $\dmodel=768$, 12 cabezas de atención y contexto de 1024 tokens. Se escribe $B$ para el tamaño del lote, $T$ para la longitud de la secuencia y $h$ para el número de cabezas. Las capas lineales se escriben aquí como en el código, $y=xW$ con $W$ de forma $(d_{\text{in}},d_{\text{out}})$.

\begin{table}[h]
\centering\small
\caption{Formas de los tensores de un Transformer decodificador.}\label{tab:formas}
\begin{tabular}{@{}lll@{}}
\toprule
Objeto & Forma & GPT-2 small, $B=8$ \\
\midrule
Identificadores de tokens & $(B,T)$ & $(8,1024)$ \\
Matriz de embeddings $E$ & $(V,\dmodel)$ & $(50\,257,768)$ \\
Activaciones & $(B,T,\dmodel)$ & $(8,1024,768)$ \\
Pesos de la primera capa del MLP & $(\dmodel,4\dmodel)$ & $(768,3072)$ \\
Puntuaciones de atención & $(B,h,T,T)$ & $(8,12,1024,1024)$ \\
Logits & $(B,T,V)$ & $(8,1024,50\,257)$ \\
\bottomrule
\end{tabular}
\end{table}

\addtocounter{figure}{1}%

\subsection*{La matriz de embeddings y el cuello de botella de la softmax}

La matriz de embeddings $E\in\R^{V\times\dmodel}$ tiene muchas más filas que columnas. En GPT-2 small, $50\,257\cdot 768=38\,597\,376$ parámetros, cerca del 31\,\% de los $124$~millones del modelo; en Llama~2 7B~\cite{llama2}, $32\,000\cdot 4096=131\,072\,000$. GPT-2 \emph{ata} los pesos de entrada y salida: los logits del token siguiente se calculan como $z=hE\T$, donde $h\in\R^{\dmodel}$ es la última activación.

\begin{proposicion}[Cuello de botella de la softmax]\label{prop:cuello}
Sean $N$ contextos con activaciones finales $h_1,\dots,h_N\in\R^{d}$, apiladas en $H\in\R^{N\times d}$, y sea $W\in\R^{V\times d}$ la matriz de salida. La matriz de logits $Z=HW\T\in\R^{N\times V}$ tiene rango a lo sumo $d$, y la matriz de log-probabilidades $L$, con $L_{cj}=\ln\softmax(z_c)_j$, tiene rango a lo sumo $d+1$.
\end{proposicion}

Yang \emph{et al.}~\cite{yang-softmax} plantean el modelado del lenguaje como una factorización de la matriz de log-probabilidades verdaderas. Si esa matriz, sobre todos los contextos posibles, tiene rango muy superior a $d+1$, ningún modelo con salida softmax de dimensión $d$ puede reproducirla exactamente. Con $V=50\,257$ y $d=768$, la restricción es severa en principio; los autores proponen una mezcla de varias softmax para superarla.

